\section{3교시 (90분) — 변경 관리: 변경은 사건이 아니라 절차다}

% =====================================================================
% 3교시 — 변경 관리 (교재 제3장, 워크북 §1)
% =====================================================================
\begin{frame}{3교시 — 이 교시에서 배우는 것}
\begin{itemize}
\item 변경의 원천 4종과 변경요청서의 3단 구조(요청서 \ra{} 영향 분석 \ra{} 의결서)
\item \textbf{6축 영향 분석} — FP 하나로는 변경의 크기를 잴 수 없다
\item 처리 경로 6규칙 — 핵심은 ② 헌장 문언 변경은 허용범위 내라도 CCB
\item 기준선 갱신 BL-02 vs 재기준선 — "반영한 뒤 편차가 줄어드는가"
\item 변경 대가 55만 원/FP · 감액도 변경 · 범위 크립의 통제 도구 셋
\end{itemize}
\keymsg{온담 변경 4건(CR-08\til{}11)이 같은 문을 통과해 BL-02가 되는 과정을 따라간다 — 4교시 실습 ①의 이론}
\note{교재 제3장 전체, 워크북 §1. Day2 2.4절이 범위 팽창의 세 형태와 처방(FP 기준선 확인서·변경 대가 경로)을 세웠다면, 이 교시는 그 경로를 실제로 통과시키는 절차를 다룬다. 이 교시의 숫자는 모두 워크북 §1.5 완성 예시본에서 가져왔고 4교시 실습 ①에서 수강생이 같은 숫자를 손으로 다시 쓴다. 강의 순서는 원천과 구조(3.1) \ra{} 6축(3.2) \ra{} 경로(3.3) \ra{} BL-02와 재기준선(3.4) \ra{} 대가와 크립(3.5) \ra{} 온담 4건과 CCB 제10차(3.6) \ra{} 핵심 정리다. 예상 소요 3분.}
\end{frame}

\begin{frame}{변경은 사건이 아니라 절차다 — 원천 4종}
\begin{tbl}[\scriptsize]{L{2.1cm}L{4.6cm}L{4.2cm}}
원천 & 뜻 & 온담 2026-09-30 \\
\midrule
요구의 변경 & 새 것을 원하거나 있던 것을 바꾼다 & CR-09 22:00 마감(서정필) \\
오류의 정정 & 기준선이 틀렸음을 발견 & CR-08 임시 연동 누락(WBS 1.1.7 "FP 0") \\
외부 사건 & 규제·벤더·타 프로젝트가 바뀜 & CR-01 배달3사 규격 v2 · CR-11 보존 요건 \\
거버넌스 변경 & 규칙·권한·정의 자체가 바뀜 & CR-10 원가 허용범위 정의(FP 0) \\
\end{tbl}
\keymsg{"요구가 바뀌었다"는 사건이다 — 그것이 기준선을 바꾸는가는 절차가 결정한다}
\note{교재 3.1. 원천이 다르면 영향 분석의 무게가 다르다 — 요구 변경은 대안 3개가 핵심이고, 오류 정정은 "왜 몰랐는가"가 등록부(R-06 OD-POS 지식)로 이어지며, 외부 사건은 귀책 분해(5교시)와 붙고, 거버넌스 변경은 FP가 0이어도 CCB다. 절차가 없으면 사건이 곧 기준선 변경이 되고 그것이 범위 크립의 정의다. 2026-09-30 시점의 변경 4건이 네 원천을 하나씩 대표한다는 점을 강조한다 — 시나리오가 그렇게 설계되었다. 실무 오류 — 요청자 칸에 "현업"·"PM"이 적히는 것(PM은 접수자다), 구두 요청을 접수하지 않는 것, 모든 변경이 "긴급"인 것. 예상 소요 6분.}
\end{frame}

\begin{frame}{변경요청서의 구조 — 세 문서가 순서대로 만들어진다}
\begin{itemize}
\item \textbf{요청서}(항목 1\til{}3) — 식별 · 내용과 사유 · 분류와 긴급도
\item \textbf{영향 분석}(항목 4\til{}6) — 6축: 범위·일정·원가·품질·리스크·계약/편익
\item \textbf{의결서}(항목 7\til{}9) — 판정과 경로 · CCB 의결·조건·반대의견 · BL-0n
\item 결과는 \textbf{변경 로그} 한 줄(항목 10) — CR-01\til{}11 누적
\item 강조 칸 둘 — "현재 기준선 문장" 인용 · 대안 3개(요청대로/축소/이연·파라미터화/기각)
\end{itemize}
\keymsg{인용이 없는 변경은 무엇에서 무엇으로 바뀌는지 모르는 변경이다 — "원래 그랬다"와 "바뀐 것이다"가 충돌한다}
\note{교재 3.1, 워크북 §1.3 항목 1\til{}10. 요청서·영향 분석·의결서는 작성 시점도 작성자도 다르다 — 요청서는 접수 시, 영향 분석은 PMO·수행사 FP 산정 후, 의결서는 CCB 후. 그래서 한 양식이 아니라 세 문서가 이어 붙은 구조로 가르친다. 워크북 예시 CR-08의 항목 2를 보이며 "현재: 범위 기술서 v1.0 §2 ① … 헌장 §5 포함 범위에 임시 연동 없음 \ra{} 요청: … FP +90"의 인용 습관과, 대안 (a)\til{}(d) 넷 중 (b) A12 편입이 채택된 이유(코드 프리즈 위반 회피)를 읽는다. 긴급도의 근거는 날짜다 — CR-08 A12 착수 12-07, CR-09 FZ-10 10-28, CR-10 10-05 성과보고, CR-11 매핑 종료 10-23. 예상 소요 6분.}
\end{frame}

\begin{frame}{6축 영향 분석 — FP 하나로 변경의 크기를 잴 수 없다}
\begin{tbl}[\scriptsize]{L{1.7cm}L{4.6cm}L{4.6cm}}
축 & 무엇을 적는가 & 온담 CR-08 (+90 FP) \\
\midrule
범위 & WBS·RTM 델타, FP, \textbf{누적} FP & 1.1.7 산출물 +3, StRS-095, 누적 +210 \\
일정 & 활동·기간·\textbf{TF}, 임계경로 여부 & A12 20\ra{}25일, TF 10\ra{}5, 임계 불변 \\
원가 & 금액 + \textbf{자금원} & +0.495억, 컨틴전시 미배정 풀 \\
품질 & TC·커버리지·밀도 분모 & TC +12, 분모 12,400\ra{}12,490 \\
리스크 & 신규·변동 EMV, 허용범위 여유 & R-19 신규 0.60, 잔여 7.26(여유 0.74) \\
계약/편익 & 변경 대가·서면·B-ID·타 프로젝트 & 증액 0.495, 합의서 6호, P3 협의 09-22 \\
\end{tbl}
\keymsg{일정 축은 TF로, 원가 축은 자금원과 함께, 계약/편익 축은 타 프로젝트·운영 영향 칸과 함께 적는다}
\note{교재 3.2, 워크북 §1.3 항목 4\til{}6. 6축은 PRINCE2 Issues practice의 6개 성과 목표(범위·일정·원가·품질·리스크·편익)에서 왔고 온담은 편익 축에 계약을 붙였다 — 고정가 계약에서 변경은 곧 대가이기 때문이다. FP 하나로 재면 CR-10은 0이고 CR-11은 마이너스라 "영향 없음"으로 읽히지만, CR-10은 예외 보고 트리거를 바꾸고 CR-11은 계약 2건과 보존 요건을 건드린다. 각 축의 값은 Day2 기준선에서 인용한다 — A12 TF 10은 §3.5 CPM 표, 미배정 풀 5.60은 §4.5 항목 6, 잔여 EMV는 §5.5. 리스크 축의 "여유 0.74 — 경고"는 허용범위 8억에 7.26이 닿아가고 있다는 5교시의 복선이다. 예상 소요 8분.}
\end{frame}

\begin{frame}{변경 4건의 6축 히트맵 — 같은 자로 다른 변경을 잰다}
\fig[0.6]{../그림/PM_Day3/fig_3_2_six_axis_radar.pdf}
\figcap{그림 3-2. CR-08\til{}11의 6축 영향 히트맵과 처리 경로. CR-10은 FP 0인데 거버넌스 축이 진하고, CR-11은 FP $-$60인데 계약 축이 진하다}
\keymsg{FP 열만 보면 CR-10은 "영향 없음"이고 CR-11은 "좋은 소식"이다 — 6축이 그 착시를 걷어낸다}
\note{교재 3.2 그림 3-2. 네 건을 한 화면에 놓고 축별로 읽는다 — 범위 축은 CR-08이 가장 크고(+90), 일정 축은 CR-09가 임계경로 A10(이미 $-$10)에 +1일을 얹어 단독으로는 허용범위 내지만 ⑫ 회복 계획 전제이며, 계약/편익 축은 CR-09가 물류본부·3PL을 건드려 운영위 보고를 부른다. 히트맵 오른쪽 처리 경로 열이 다음 프레임의 규칙 6개로 이어진다. 수강생에게 "여러분 회사 변경 양식에 이 여섯 열 중 몇 개가 있는가"를 묻는다 — 대개 범위·원가 둘이다. 예상 소요 5분.}
\end{frame}

\begin{frame}{처리 경로 6규칙 — 허용범위는 크기의 기준, 헌장은 권한의 기준}
\begin{tbl}[\scriptsize]{L{0.5cm}L{5.6cm}L{4.8cm}}
\# & 조건 & 경로 \\
\midrule
① & 6축 모두 허용범위 내 · 헌장 §5·§12·§13 문언 불변 · 타 영향 없음 & \textbf{P1 PM 전결}(RACI 16행), 다음 CCB 사후 보고 \\
② & \textbf{헌장 문언 변경} 또는 어느 축 초과 & \textbf{CCB} 승인(17행); 수행사 FP 영향 시 \textbf{변경협의체} 합의 선행(18행) \\
③ & 관리예비비 필요 & \textbf{스폰서} 결재(20행) \\
④ & 타 프로젝트·운영 조직 영향 & 프로그램 \textbf{운영위} 보고 병행 \\
⑤ & 릴리스 이연(R5) & 헌장 §13 인용, CCB \\
⑥ & 승인된 변경 & CCB 후 5영업일 내 \textbf{BL 번호} 일괄 갱신 \\
\end{tbl}
\keymsg{CR-08(+90)·CR-11($-$60)은 허용범위 안이다 — 그러나 헌장 §5 문장을 바꾸므로 CCB다(규칙 ②)}
\note{교재 3.3, 워크북 §1.5 첫머리 "처리 경로 규칙(변경 관리 계획 v1.0 §3, G1 확정 [추가 설정])". 두 문이 다르다는 것이 이 교시의 핵심 문장이다 — 허용범위는 "얼마나 큰가"를 재고, 헌장은 "누가 서명했는가"를 잰다. 스폰서와 이사회가 서명한 문장을 PM이 전결로 바꾸면 다음 게이트에서 헌장과 기준선이 어긋난다. 순서도 중요하다 — 변경협의체(계약, 수행사와 FP 합의) \ra{} CCB(사내 의결) \ra{} 계약 서면. 변경협의체 없이 CCB가 FP를 정하면 수행사는 G4에서 다시 협상한다. 민간에는 공공의 과업심의위원회가 없으므로 이 협의체가 그 자리를 대신한다는 점, CR-07(+40, 헌장 문언 불변)은 규칙 ①로 전결·사후 보고된 대조 사례라는 점을 짚는다. 예상 소요 8분.}
\end{frame}

\begin{frame}{네 건이 같은 문을 통과한다 — 접수에서 BL-02까지}
\fig[0.6]{../그림/PM_Day3/fig_3_6_change_flow.pdf}
\figcap{그림 3-6. 접수 \ra{} 6축 \ra{} 경로 판정 \ra{} 변경협의체(09-30) \ra{} 운영위(10-08) \ra{} CCB(10-15) \ra{} 계약 서면(10-23) \ra{} BL-02(10-23). CR-07 전결은 대조}
\keymsg{변경협의체 \ra{} CCB \ra{} 계약 서면 \ra{} BL — 순서가 바뀌면 G4에서 대가를 다시 협상한다}
\note{교재 3.6 그림 3-6. 네 건의 날짜를 따라 읽는다 — 접수 09-16\til{}22(모두 FZ-10 접수 마감 10-14 전), 변경협의체 09-30 합의(CR-08·09·11 FP 확인, 이재현), 운영위 10-08(CR-09 물류·3PL 영향), CCB 제10차 10-15, 계약 서면 10-23(SI 변경 합의서 6호에 세 건 통합 +65 FP·+0.3575억, 용역 변경 계약 1호 $-$2.0), BL-02 발행 10-23. CR-07은 같은 그림의 옆길 — PM 전결 09-17, CCB 사후 보고. 규칙 ⑥ 5영업일이 10-15 \ra{} 10-23으로 지켜졌음을 확인시킨다. 예상 소요 5분.}
\end{frame}

\begin{frame}{CR-08 · CR-09 — 범위 추가와 요구 변경은 어떻게 다르게 읽히는가}
\begin{columns}[T]
\begin{column}{0.49\textwidth}
\subhead{CR-08 파일럿 임시 재고 동기화 (+90 FP)}
\begin{itemize}
\item 헌장 v1.1 수정 요청 1번 · 정미란 · 09-16
\item 대안 (b) A12 편입 — 코드 프리즈 회피
\item +0.495억 미배정 풀 · TF 10\ra{}5 · R-19 신규
\item 헌장 §5 문언 변경 \ra{} 변경협의체 \ra{} \textbf{CCB}
\end{itemize}
\end{column}
\begin{column}{0.49\textwidth}
\subhead{CR-09 가맹 발주 마감 20:00\ra{}22:00 (+35 FP)}
\begin{itemize}
\item 협의회 공문 · 서정필(S-10) · 09-17
\item 대안 (c) \textbf{파라미터화}, 초기 운용 20:00
\item +0.19억 · 임계 A10 +1일 · R-16 P 30\ra{}45
\item 물류·3PL 영향 \ra{} 규칙 ④ \textbf{운영위} + CCB
\end{itemize}
\end{column}
\end{columns}
\keymsg{CR-09가 CCB인 이유는 크기가 아니다 — P1 CCB는 이천센터 야간조 편성을 결정할 수 없다}
\note{워크북 §1.5 CR-08·CR-09 전문, 교재 3.6·해설 "왜 CR-09를 대안 (c)로". CR-08은 오류 정정에 가까운 범위 추가다 — WBS 1.1.7 카드가 "FP 0 — 배포 기록만"이었고, 2019년 스마트오더가 매장 운영 부하로 중단된 경로를 12개점에서 재현하지 않으려면 연동 2본(68 FP)과 배포·회수 도구(22 FP)가 필요하다. CR-09는 요구 자체는 타당하나(현행 웹 23:00보다 후퇴), 22:00 마감은 P1 화면이 아니라 21:00 피킹·폴백 야간 배치 21:00·03:00·대성 배송 출발 시각을 건드린다. 그래서 파라미터로 구현해 두고 운용 전환은 야간조 편성안 첨부 후 2027-01 CCB 재상정으로 조건을 달았다. 한동석(물류본부)은 CCB 위원이 아니고 파라미터 구현에 이의도 없었지만 서면 의견을 "반대의견" 칸에 실명 기재했다 — Day1 거버넌스의 CCB 규칙이자, 재상정 때 그의 의견이 출발점이 되게 하려는 장치다. 예상 소요 8분.}
\end{frame}

\begin{frame}{CR-10 · CR-11 — FP 0과 FP $-$60에도 6축을 다 돌린다}
\begin{columns}[T]
\begin{column}{0.49\textwidth}
\subhead{CR-10 원가 허용범위 정의 확정 (FP 0)}
\begin{itemize}
\item 헌장 §12 "+6\% (10.08)" \ra{} "컨틴전시 전액 10.2 (+7.0\%)"
\item 세 문서가 두 값 — 155.88 vs 156.0
\item 예외 보고 트리거 "EAC $>$ 156.0" 단일화
\item \textbf{CCB} + 스폰서 확인 + 운영위 보고
\end{itemize}
\end{column}
\begin{column}{0.49\textwidth}
\subhead{CR-11 전환 대상 5년 한정 ($-$60 FP)}
\begin{itemize}
\item 헌장 §5 ⑤ 7.4TB \ra{} 5년 이내 약 4.8TB
\item 재경팀장 보존 요건 확인서 09-22
\item SI \textbf{감액 $-$0.33억} · 용역 계약 11.4\ra{}9.4
\item A09 TF 5\ra{}15 · 계약 서면 2건 · CCB
\end{itemize}
\end{column}
\end{columns}
\keymsg{CR-10은 돈이 한 푼도 움직이지 않지만 예외 보고의 자를 바꾼다 — 감액 변경도 변경이다}
\note{워크북 §1.5 CR-10·CR-11 전문, 교재 해설 "왜 CR-11에도 6축을 다 돌렸는가". CR-10은 Day2 원가 기준선 항목 9에서 발견된 불일치의 확정이다 — 프로그램 규정 별표의 +6\%는 승인 예산 168.0을 분모로 한 프로그램 계층 규정이고, P1의 PM 허용범위는 "스폰서 보고 없이 쓸 수 있는 돈" = 컨틴전시 10.2다. 10.08로 두면 계획 원가 + 허용범위 = 155.88 $\neq$ 집행 한도 156.0이라 "허용범위 내인데 집행 한도 초과"가 논리적으로 가능해진다. 10-05 성과보고의 첫 EAC 기반 예외 판정이 이 정의를 쓰므로 긴급도가 있다. CR-11의 핵심은 원가가 아니라 보존 요건이다 — 확인서 없이 축소하면 세무조사 때 "5년 전 거래를 왜 못 보느냐"가 P1 책임이 된다. 원가 기준선 BL-01은 이미 9.4로 기반영되어 있었고 용역 계약만 11.4였으므로 이 변경은 계약을 기준선에 정합시키는 것이다. 관리예비비 복귀 조건 1건이 해제된다. 예상 소요 7분.}
\end{frame}

\begin{frame}{변경 로그 CR-01\til{}CR-11 — 누적 +185FP는 허용범위 ±372의 49.7\%}
\begin{tbl}[\scriptsize]{L{1.2cm}L{4.2cm}rrL{2.6cm}}
CR & 내용 & FP & 억 & 경로 \\ \midrule
03 & 인터페이스 코드 체계 보정 & +22 & +0.12 & PM 전결 \\
06 & 가맹 포털 알림 채널 추가 & +58 & +0.32 & CCB \\
07 & 인터페이스 미등재 2본 실체화(IS-03) & +40 & +0.22 & CCB \\
08 & 파일럿 임시 연동(헌장 §5 문언) & +90 & +0.495 & CCB · 조건 4 \\
09 & 22:00 마감 파라미터(조건부) & +35 & +0.19 & 협의체 \ra{} CCB \\
11 & 데이터 전환 5년 한정(감액) & $-$60 & $-$0.33 & CCB \\
\end{tbl}
\keymsg{감액도 변경이다 — CR-11이 로그에 없으면 "범위가 늘기만 했다"는 착시가 남고 변경 대가 협상에서 60FP를 잃는다}
\note{교재 3.6 "변경 로그", 워크북 §1.5 변경 로그 11건과 §5.5 변경 대가 산정. CR-01·02·04·05·10은 FP 0(절차·문언·거버넌스 변경)이라 표에서 뺐다 — 그러나 로그에는 있다는 것이 요점(변경 = FP가 아니라 기준선을 건드리는 모든 것). 누적 +185FP = 1.5\%, 허용범위 ±3\% = ±372의 49.7\%. 금액 +1.0175억(SI, 55만 원/FP)이 BL-02에서 PMB로 편입되고 미배정 풀 2.47이 그만큼 줄었다(3교시 앞 프레임 "네 건이 같은 문을 통과한다"의 결과). 학생에게 묻기 — "허용범위의 절반을 9개월에 썼다. 남은 8개월은?" 답은 정답이 없고, R-10(정본 판정 후 변경 요청)의 EMV 0.76이 그 질문의 정량판이다. 예상 소요 4분.}
\end{frame}

\begin{frame}{"Day1 문서는 고쳐진다" — 헌장 v1.1과 변경의 흐름}
\fig[0.48]{../그림/PM_Day3/fig_3_6_change_flow.pdf}
\figcap{그림 3-6. 접수 \ra{} 분류 \ra{} 6축 영향 \ra{} 경로 판정 \ra{} 의결 \ra{} 기준선 갱신 BL-0n \ra{} 문서 역전파(헌장·RTM·WBS 사전·등록부·PV)}
\keymsg{헌장을 고치는 것은 실패가 아니다 — 고치지 않은 헌장이 실패다. 단, 고친 흔적(v1.1, CR 번호, 승인자)이 없으면 그것은 변경이 아니라 변조다}
\note{교재 3.7, 워크북 §1.6 해설과 §6. Day2 §6이 남긴 헌장 v1.1 수정 요청 10건 중 오늘 CCB에 오른 것이 CR-08·10·11(3건)이고, 나머지는 BL-02 발행 시 문서 갱신으로 흡수됐다. 역전파 순서 — 헌장 §5 문언(CR-08) \ra{} 범위 기술서·WBS v1.1 \ra{} RTM v1.2 \ra{} PV(2026-12 +0.33·2027-01 +0.69) \ra{} 등록부 v1.2. 이 순서가 거꾸로 되면(PV부터 고치고 헌장은 나중에) 감리가 "기준선 근거 없음"으로 잡는다 — 감리 1번 지적이 정확히 그것이었다. Day1 슬라이드 "헌장은 권한 문서다"를 다시 띄우고 "권한 문서를 누가 고칠 수 있는가"로 닫는다(스폰서 서명 = 정미란). 예상 소요 4분.}
\end{frame}

\begin{frame}{3교시 핵심 정리}
\begin{enumerate}
\item 변경은 사건이 아니라 절차 — 접수·6축·경로·의결·갱신의 다섯 단계 중 하나라도 건너뛰면 변조다
\item 6축 영향 분석은 FP 하나로 대체되지 않는다 — CR-10은 FP 0인데 원가 축 정의를 바꿨다
\item 처리 경로 6규칙 — 허용범위는 크기의 기준, 헌장 문언은 권한의 기준, 운영 영향은 소관의 기준
\item BL-02: FP 12,585 · PMB 146.82 · 집행 한도 156.00 불변 — 기준선 갱신은 승인된 변경으로만
\item 로그 누적 +185FP(49.7\%)와 감액 $-$60 — 변경 로그는 협상의 근거이자 허용범위의 잔고
\end{enumerate}
\keymsg{4교시 실습 — CR-09를 직접 6축으로 재고 의결서를 쓴 뒤, 9/30 세 숫자에서 열두 값을 손으로 만든다}
\note{교재 제3장 핵심 정리 1\til{}5. 생각해 볼 질문(제3장 3 — "여러분 회사에서 마지막으로 헌장이 개정된 것은 언제인가, 개정된 적이 없다면 그것은 무엇을 뜻하는가"). 점심 전 마지막 이론 교시이므로 4교시 준비물(워크북 §1.4·§2.4 빈 템플릿, 계산기)을 안내한다. 예상 소요 3분.}
\end{frame}
